
\begin{cnabstract}
2012年7月，欧洲核子研究中心大型强子对撞机上的ATLAS和CMS实验组宣布发现Higgs玻色子，标准模型的粒子谱
由此得到补全。实验表明，新发现粒子的性质与标准模型Higgs玻色子的预言大致相符。
为了确认该粒子是标准模型的Higgs玻色子，还是新物理模型中的Higgs粒子，并进一步理解电弱自发对称破缺机制，
需要更精确的实验检验。为此，首先需要对Higgs粒子各产生道的信号和背景给出精确的理论预言。
此外，检验矢量玻色子的规范耦合、研究Higgs粒子的自耦合及Higgs势也同样重要。
尽管标准模型取得了巨大成功，但它并非终极理论，因此仍需探索新物理规律。

为了更精确地检验W玻色子对与Z玻色子或光子的耦合，我们研究了LHC上
$pp\to W^+W^-+{\rm jet}+X$过程的次领头阶QCD和电弱修正，并考虑了W玻色子衰变的自旋关联效应。
在LHC Run 2阶段，更高的对撞能量和亮度使得产生横动量达到TeV量级的规范玻色子成为可能；
此时，次领头阶电弱修正因Sudakov效应而尤为重要。2013年的Les Houches国际高能物理会议也强调，
亟需计算矢量玻色子对伴随jet产生过程及玻色子衰变的次领头阶电弱和QCD修正。
我们研究了领头阶及次领头阶截面对领头jet横动量截断值的依赖关系，并采用MadSpin方法处理W玻色子衰变的
自旋关联效应。我们还给出了末态粒子的横动量和快度分布；结果表明，在高横动量区域，电弱修正的贡献十分显著。

为了研究Higgs粒子的自耦合性质，我们计算了双Higgs二重态模型类型II下矢量玻色子聚合过程
$pp\to h^0h^0+2jets+X$的次次领头阶QCD修正。
在标准模型下，矢量玻色子聚合产生Higgs粒子对有着第二大的散射截面。
此外，该过程具有两个位于前向和后向区域的jet这一清晰的实验特征，可用于研究Higgs粒子的自耦合并重构Higgs势。
我们采用结构函数法计算了此过程的次次领头阶QCD修正，研究了总截面对模型参数的依赖、因子化和重整化能标的
不确定度，以及PDF和$\alpha_s$的不确定度。结果表明，次次领头阶QCD修正可以显著减小能标不确定度。
我们还给出了末态Higgs粒子的运动学分布，可用于研究Higgs粒子的自耦合
（$\lambda_{H^0h^0h^0}$和$\lambda_{h^0h^0h^0}$）。


本论文的创新点如下：
\begin{itemize}
\item 我们首次计算了LHC上$pp\to W^+W^-+{\rm jet}+X$过程的次领头阶电弱修正，
给出了次领头阶QCD与电弱修正的组合结果，并考虑了W玻色子衰变的自旋关联效应。
计算中我们巧妙地利用CKM 矩阵的幺正性极大地简化了计算量。另外，通过改变领头jet的横动量截断值，
我们验证了电弱Sudakov效应。
在处理W玻色子衰变的时候，我们用MadSpin方法来包含W玻色子衰变的自旋关联效应。

\item 在处理过程$pp\to WW+{\rm jet}+X$的单圈张量积分时，会遇到由小Gram行列式带来的数值计算不稳定问题，
我们开发了双精度和四精度联合计算的程序包，解决了数值不稳定问题，并显著提高了计算效率。
对于实辐射过程涉及到的$W^+W^-+jet+\gamma$四体事例，我们采用精确的事例选择规则来区分jet和光子事例，
并利用夸克-光子碎裂函数处理红外发散问题。

\item 我们首次采用结构函数法计算了双Higgs二重态模型类型II下矢量玻色子聚合过程
$pp\to h^0h^0+2jets+X$的次次领头阶QCD修正。结果显示次次领头阶的计算使得能标不确定度极大地减小。
此外，我们开发的程序包高效灵活，能计算矢量玻色子聚合道产生单个或多个Higgs粒子的过程，
适用于标准模型和双Higgs二重态模型。
\end{itemize}

\keywords{标准模型，双Higgs二重态模型，QCD修正，电弱修正，次次领头阶}
\end{cnabstract}


\begin{enabstract}
In July 2012, the ATLAS and CMS collaborations at the CERN Large Hadron Collider (LHC)
announced the discovery of a Higgs boson, completing the particle spectrum of the standard model (SM).
Measurements indicate that the properties of this particle are broadly consistent with those predicted
for the SM Higgs boson. More precise tests are needed to establish whether it is the SM Higgs boson
or a Higgs particle associated with new physics, and to deepen our understanding of electroweak
symmetry breaking (EWSB). Such tests require precise theoretical predictions for the signals and
backgrounds in Higgs production channels. Measurements of vector-boson gauge couplings, studies of
Higgs self-couplings and the Higgs potential, and searches for new physics are also important.

To test the couplings of W bosons more precisely, we investigate the next-to-leading order (NLO)
electroweak (EW) and QCD corrections to the process
$pp\to W^+W^-+{\rm jet}+X \to e^+ \mu^- \nu_e \bar \nu_{\mu}+{\rm jet}+X$ at the LHC.
The higher collision energy and luminosity of LHC Run 2 make it possible to produce gauge bosons
with transverse momenta at the TeV scale, where NLO EW corrections become important because of
the Sudakov effect. The 2013 Les Houches workshop also emphasized the need for NLO EW and QCD
corrections to vector-boson pair production in association with a jet, including boson decays.
We study the dependence of the leading-order (LO) and NLO cross sections on the leading-jet
transverse-momentum cut. We use MadSpin to include spin correlations in the W-boson decays and
present the transverse-momentum and rapidity distributions of the final-state particles.
The results show that EW corrections are significant at high transverse momentum.

To study Higgs self-couplings, we calculate the next-to-next-to-leading order (NNLO) QCD
corrections to the vector-boson fusion (VBF) process $pp\to h^0h^0+2jets+X$ in the type-II
two-Higgs-doublet model (2HDM). In the SM, VBF has the second-largest cross section among
Higgs-pair production channels. Its characteristic final state contains two Higgs bosons and
two jets in the forward and backward rapidity regions, allowing the Higgs self-couplings
and potential to be studied. Using the structure-function approach, we examine the dependence
of the total cross section on model parameters and evaluate uncertainties associated with the
factorization and renormalization scales, parton distribution functions (PDFs), and $\alpha_s$.
We find that NNLO QCD corrections significantly reduce the scale uncertainty. We also present
kinematic distributions of the final-state Higgs bosons, which can be used to study the
self-couplings $\lambda_{H^0h^0h^0}$ and $\lambda_{h^0h^0h^0}$.


The main contributions of this thesis are as follows:
\begin{itemize}
\item We calculate the NLO EW corrections to the process
$pp\to W^+W^-+{\rm jet}+X \to e^+ \mu^- \nu_e \bar \nu_{\mu}+{\rm jet}+X$
at the LHC and combine them with the NLO QCD corrections. We use the unitarity of the CKM matrix
to simplify the calculation and study the EW Sudakov effect by varying the leading-jet
transverse-momentum cut. MadSpin is used to include spin correlations in the W-boson decays.

\item Small Gram determinants cause numerical instabilities in the one-loop tensor integrals
for $pp\to WW+jet+X$. We developed a program that combines double and quadruple precision
to address these instabilities and improve computational efficiency. For the real-photon
radiation contribution, we use a precise event-selection rule to distinguish jets from
photons in $W^+W^-+jet+\gamma$ final states and a quark-to-photon fragmentation function
to handle the associated infrared divergence.

\item 
Using the structure-function approach, we calculate the NNLO QCD corrections to the VBF
process $pp\to h^0h^0+2jets+X$ in the type-II 2HDM. These corrections significantly reduce
the scale uncertainty. Our code can also calculate the production of one or more Higgs
bosons via VBF in both the SM and the 2HDM.
\end{itemize}

\enkeywords{Standard Model, Two Higgs Doublet Model, QCD correction, Electroweak correction, 
next-to-next-to-leading order}

\end{enabstract}
